\chapter{Build system}
\label{chap:build}
\sloppy

\PR is built with the GNU autotools. This chapter describes how the
\code{configure} script chooses the compiler flags and the preprocessor
definitions, and how to change it. The user side (prerequisites, options,
\code{make install}) is in the User Guide.

\section{Files}

\begin{longtable}{L{3cm} L{10.5cm}}
  \toprule
  \textbf{File} & \textbf{Content} \\
  \midrule
  \endhead
  \bottomrule
  \endfoot
  \code{configure.ac} & Options, compiler detection and flags, library
    checks, the substituted variables and the list of generated
    \code{Makefile}s. \\
  \code{acinclude.m4} & The \PR macros: NetCDF and Parallel NetCDF
    discovery (\code{AX\_PROG\_NC\_CONFIG}, \code{AX\_PROG\_NF\_CONFIG},
    \code{RR\_PATH\_NETCDF}, \code{RR\_PATH\_NETCDF\_F90},
    \code{RR\_NETCDF4}, \code{RR\_CDF5}, \code{AX\_PROG\_PNETCDF\_CONFIG},
    \code{RR\_PATH\_PNETCDF}), the MPI compiler (\code{ACX\_MPI},
    \code{RCM\_MPI\_CHECK\_MPI3}), Fortran feature checks
    (\code{RCM\_FC\_CHECK\_IEEE\_ARITHMETIC},
    \code{RCM\_FC\_CHECK\_QUAD\_PRECISION}) and the compiler identification
    (\code{RCM\_FC\_VENDOR}). \\
  \code{bootstrap.sh} & Runs \code{autoreconf -f -i}. \\
  \code{makeinc} & Included by every \code{Makefile.am}: the
    \code{GITVER} definition (\code{-DGIT\_VER} from \code{git describe} or
    the \code{tag} file) and the CLM~4.5 scheme switches. \\
  \code{*/Makefile.am} & One per directory, with explicit compile rules
    and module dependencies. \\
\end{longtable}

The generated files (\code{configure}, \code{Makefile.in}, \code{aclocal.m4},
\code{config.h}, \ldots) are not in the repository. \code{config.h} is
written by \code{AC\_CONFIG\_HEADERS} but no source file includes it: all
the definitions reach the code as \code{-D} flags.

\section{Order of the checks}

\code{configure.ac} runs in this order:

\begin{enumerate}
  \item The \code{--with} and \code{--enable} options. Each option that
    selects code adds its definition to \code{AM\_CPPFLAGS}.
  \item The NetCDF (and Parallel NetCDF) paths, through \code{nc-config}
    and \code{nf-config} when available.
  \item The Fortran compiler vendor (Section~\ref{sec:fcvendor}) and the
    \code{COMPILER\_*} conditionals.
  \item The base flags of the vendor, then the debug or production flags,
    then the \code{--enable-reproducible} and OpenACC flags
    (Section~\ref{sec:fcflags}).
  \item The NetCDF features (NetCDF4/HDF5, filters, CDF5), the Fortran
    feature checks, which add \code{-DF2008} and \code{-DQUAD\_PRECISION}
    to \code{FCFLAGS} when supported, and the MPI compiler with its MPI-3
    check. \code{configure} stops if the MPI library lacks the MPI-3
    interfaces.
  \item The check that \code{MPIFC} and \code{FC} are the same compiler,
    the OASIS libraries, the final \code{CPPFLAGS} and \code{LIBS}.
  \item One list of \code{AC\_SUBST} calls and \code{AC\_OUTPUT}.
\end{enumerate}

\code{AC\_SUBST} only marks a variable for substitution: the value written
to the \code{Makefile}s is the one the variable has at \code{AC\_OUTPUT}.
The substitutions are therefore listed once, at the end, and a variable
can be changed anywhere before.

\section{Options}

Every \code{--enable} option follows the same pattern: a variable with
default \code{0} (off), set to \code{1} only when the option is given with a
value other than \code{no}, so that \code{--disable-OPTION} and
\code{--enable-OPTION=no} leave the option off:

\begin{lstlisting}
USE_FOO="0"
AC_ARG_ENABLE(foo,
  AS_HELP_STRING([--enable-foo], [Compile the foo option.]),
  [AS_IF([test "x$enableval" != xno], [USE_FOO="1"])]
)
if test "x$USE_FOO" = x1; then
  AM_CPPFLAGS="-DFOO $AM_CPPFLAGS"
fi
\end{lstlisting}

Add a \code{-D} definition only if the code reads it, and an
\code{AM\_CONDITIONAL} only if a \code{Makefile.am} needs it. When several
options set the same variable (the four OpenACC options set
\code{USE\_OPENACC} to 1--4), their actions run in the order of the
\code{AC\_ARG\_ENABLE} calls in \code{configure.ac}, not in the order of the
command line, so the last declared of the given options wins.

The preprocessor definitions set by \code{configure} are:

\begin{longtable}{L{4.5cm} L{9cm}}
  \toprule
  \textbf{Definition} & \textbf{Set by} \\
  \midrule
  \endhead
  \bottomrule
  \endfoot
  \code{CLM45} & \code{--enable-clm45} \\
  \code{OASIS}, \code{ECLM} & \code{--enable-oasis}; \code{--enable-eclm}
    sets both \\
  \code{CPL} & \code{--enable-cpl} \\
  \code{DEBUG} & \code{--enable-debug} \\
  \code{SINGLE\_PRECISION\_REAL} & \code{--enable-singleprecision} \\
  \code{OPENACC}, \code{STDPAR} & any OpenACC option; \code{STDPAR} only
    with \code{--enable-openacc-stdpar} \\
  \code{ASYNC\_NETCDF} & \code{--enable-async-netcdf} \\
  \code{PNETCDF}, \code{PNETCDF\_IN\_NETCDF}, \code{NETCDF4\_COMPRESS} &
    \code{--enable-pnetcdf}, \code{--enable-parallel-nc},
    \code{--enable-nc4-filters}; each later one removes the earlier
    definitions \\
  \code{NETCDF4\_HDF5}, \code{NCFILTERS\_AVAIL}, \code{NETCDF\_CDF5} &
    \code{RR\_NETCDF4} and \code{RR\_CDF5}, when the NetCDF library supports
    the feature \\
  \code{F2008}, \code{QUAD\_PRECISION} & the Fortran feature checks, added
    to \code{FCFLAGS} \\
  \code{IBM} & IBM XL compiler, added to \code{FCFLAGS} \\
\end{longtable}

The processor options (\code{--enable-knl}, \code{-zen3}, \code{-bdw},
\code{-nhl}, \code{-skl}) and \code{--enable-profile} define nothing: they
only add compiler flags.

\section{Compiler identification}
\label{sec:fcvendor}

\code{RCM\_FC\_VENDOR(COMMAND, VARIABLE)} identifies the compiler run by
\code{COMMAND} from the banner it prints, not from the command name, so
that versioned names (\code{gfortran-13}), MPI wrappers (\code{mpifort},
\code{mpiifx}) and the Cray \code{ftn} wrapper are recognised. Leading
words that start with \code{scorep} or \code{-} are skipped, so a Score-P
prefix is ignored. The macro runs the compiler with \code{--version},
\code{-V} and \code{-qversion} (each compiler answers at least one, the
others just fail) and matches the output, in this order:

\begin{longtable}{L{4.5cm} L{2.5cm} L{6.3cm}}
  \toprule
  \textbf{Banner contains} & \textbf{Vendor} & \textbf{Compiler} \\
  \midrule
  \endhead
  \bottomrule
  \endfoot
  \code{Cray Fortran} & \code{cray} & Cray CCE \\
  \code{NAG Fortran} & \code{nag} & NAG nagfor \\
  \code{IBM XL Fortran}, \code{IBM Open XL Fortran} & \code{ibm} &
    IBM XL \\
  \code{nvfortran}, \code{pgfortran}, \code{pgf9}$n$ & \code{nvhpc} &
    NVIDIA HPC SDK, PGI \\
  \code{(IFX)} & \code{intel\_llvm} & Intel ifx \\
  \code{(IFORT)} & \code{intel} & Intel ifort \\
  \code{GNU Fortran} & \code{gnu} & GNU gfortran \\
  anything else & \code{unknown} & no vendor flags, with a warning \\
\end{longtable}

The order matters because a wrapper banner can name more than one
compiler: Cray CCE comes first, so that \code{ftn} is \code{cray} only when
it fronts CCE and is otherwise the compiler it fronts. The macro also sets
\code{VARIABLE\_major} to the major version of ifort, which chooses between
\code{-openmp} (before version 17) and \code{-qopenmp}.

\code{configure} calls the macro twice: on \code{FC}, to choose the flags
(\code{fc\_vendor}), and on \code{MPIFC} after \code{ACX\_MPI}
(\code{mpifc\_vendor}). The model and most tools are compiled with
\code{MPIFC}, so a warning is printed when the two differ.

The vendor sets the automake conditionals \code{COMPILER\_GNU},
\code{COMPILER\_INTEL}, \code{COMPILER\_INTEL\_LLVM}, \code{COMPILER\_PGI}
(for \code{nvhpc}), \code{COMPILER\_NAG}, \code{COMPILER\_IBM} and
\code{COMPILER\_CRAY}. Only \code{COMPILER\_PGI} is used, in
\code{Main/Makefile.am} and \code{Main/mpplib/Makefile.am}, to work around
nvfortran crashes on single files (\code{-Mnoinline} or \code{NOOPTFC}).

To support a new compiler: add its banner to \code{RCM\_FC\_VENDOR} before
any vendor its banner may also mention, add a branch to the three
\code{case \$fc\_vendor} blocks (base, debug and production flags) and, if
the \code{Makefile}s need it, a conditional.

\section{Compiler flags}
\label{sec:fcflags}

\code{configure} produces three Fortran flag sets:

\begin{description}
  \item[\code{FCFLAGS}] The flags of every Fortran compilation.
  \item[\code{NOOPTFC}] The flags for files compiled with less (or no)
    optimisation: the RRTMG k-distribution tables when they are compiled
    into the code (\code{Main/radlib}) and the files worked around for
    nvfortran.
  \item[\code{OMPFLAGS}] The OpenMP flag. It is used in \code{Share}, the
    pre- and post-processing programs and the tests, not in the model
    (\code{Main}), which is parallel through MPI only.
\end{description}

The vendor branches first set the base flags: big-endian unformatted I/O,
OpenMP and the processor options. The debug (\code{--enable-debug}) or
production branch then sets the optimisation and checks. The optimisation
and check flags are put in front of the \code{FCFLAGS} given by the user,
so that the user flags come later on the command line and win. In debug,
\code{NOOPTFC} is the same as \code{FCFLAGS} for every vendor.

When the user gives no \code{FCFLAGS}, \code{AC\_PROG\_FC} sets a default
(\code{-g}, plus \code{-O2} for a compiler it takes for GNU).
\code{configure} records at its start whether \code{FCFLAGS} was given
(\code{rcm\_user\_fcflags}); the \code{gnu}, \code{intel},
\code{intel\_llvm} and \code{nvhpc} branches drop the default, and keep the
flags of the user. For \code{gnu} this matters: the default \code{-O2}
would come after, and override, the \code{-O3} of production and the
\code{-O0} of debug. The \code{cray} branch removes \code{-g} and
\code{-O2} from \code{FCFLAGS} in any case; \code{nag} and \code{ibm} keep
the default \code{-g}.

\begin{description}
  \item[\code{gnu}] \mbox{}
    \begin{itemize}
      \item Base: \code{-fconvert=big-endian -fno-range-check -fPIC};
        OpenMP \code{-fopenmp}; \code{-march=} for the processor options.
      \item Production on Linux: \code{-O3 -finline-functions
        -mtune=native -ftree-vectorize -march=native}; on macOS:
        \code{-O3 -finline-functions}. \code{NOOPTFC}: \code{-O2
        -finline-functions}.
      \item Debug: \code{-O0 -g -fbacktrace -fdump-core}, the warnings
        and \code{-fcheck=array-temps,bounds,do,mem,recursion}; on Linux
        also \code{-pg}, \code{-fcheck=pointer} and
        \code{-ffpe-trap=zero,invalid,overflow,underflow}.
    \end{itemize}
  \item[\code{intel}] \mbox{}
    \begin{itemize}
      \item Base: \code{-convert big\_endian
        -assume byterecl -heap-arrays -fPIC}; OpenMP \code{-qopenmp};
        \code{-x}/\code{-ax} for the processor options.
      \item Production: \code{-ftz -ip -O3 -finline-functions -fp-model
        source}. \code{NOOPTFC}: \code{-ftz -O0 -fp-model source}.
      \item Debug: \code{-O0 -fpe0 -g -check all -traceback -warn all -ftz
        -fimf-arch-consistency=true -fltconsistency -mieee-fp}.
    \end{itemize}
  \item[\code{intel\_llvm}] \mbox{}
    \begin{itemize}
      \item Base: as \code{intel} plus
        \code{-qopt-report}; OpenMP \code{-fiopenmp
        -fno-openmp-target-do-concurrent}, none in debug.
      \item Production: \code{-ftz -O3 -flto -finline-functions -fp-model
        source}. \code{NOOPTFC}: \code{-ftz -O0 -fp-model source}.
      \item Debug: \code{-O0 -g -traceback -warn all -check
        bounds,contiguous -mieee-fp -fpe0}.
    \end{itemize}
  \item[\code{nvhpc}] \mbox{}
    \begin{itemize}
      \item Base: \code{-byteswapio}; OpenMP
        \code{-mp}.
      \item Production: \code{-O3 -fast -fastsse}. \code{NOOPTFC}:
        \code{-O1}.
      \item Debug: \code{-Mbounds -Mcache\_align -Mchkstk -Mdalign
        -Miomutex -Mrecursive -Ktrap=divz -O0 -g -gopt -traceback}.
    \end{itemize}
  \item[\code{nag}] \mbox{}
    \begin{itemize}
      \item Base: \code{-convert=big-endian -pic -f2018}; OpenMP
        \code{-openmp}.
      \item Production: \code{-O3}. \code{NOOPTFC}: \code{-O2}.
      \item Debug: \code{-O0 -g -pg -mtrace=all
        -gline -ieee=nonstd -nan -Warn=allocation -C=all}.
    \end{itemize}
  \item[\code{ibm}] \mbox{}
    \begin{itemize}
      \item Base: \code{-DIBM}; OpenMP \code{-qsmp=omp}.
      \item Production: \code{-qsuppress=1501-245 -O3 -qstrict
        -qxlf2003=polymorphic -qmaxmem=-1 -qzerosize}. \code{NOOPTFC}:
        the base flags only.
      \item Debug: \code{-O0 -qsuppress=1501-245 -qcheck
        -qflttrap=nanq,overflow -qxlf2003=polymorphic -qmaxmem=-1
        -qzerosize}.
    \end{itemize}
  \item[\code{cray}] \mbox{}
    \begin{itemize}
      \item Base: removes the autoconf defaults \code{-g}, which limits
        the CCE optimisation, and \code{-O2}, which would override the
        level set later; then \code{-h byteswapio}; OpenMP
        \code{-fopenmp}.
      \item Production: \code{-O3}. \code{NOOPTFC}: the base flags only.
      \item Debug: \code{-O0 -g}.
    \end{itemize}
\end{description}

\code{--enable-reproducible} works only with \code{gnu}. It adds
\code{-ffp-contract=off}, so that no \code{a*b+c} is contracted to a fused
multiply-add in one part of a loop and not in the other, and
\code{-nostdinc} with an explicit \code{-fintrinsic-modules-path}: some
Linux distributions make gfortran pre-include
\code{math-vector-fortran.h}, so that vectorised loops call the glibc
vector math library, whose results differ in the last bit from the scalar
functions. With both, the results do not depend on the MPI domain
decomposition.

The OpenACC options work only with \code{nvhpc}; otherwise
\code{configure} warns and sets no GPU flag. They add to \code{FCFLAGS},
\code{NOOPTFC} and \code{LDFLAGS} (\code{-traceback
-cudalib=nvtx,cusparse,curand}):

\begin{longtable}{L{4.5cm} L{9cm}}
  \toprule
  \textbf{Option} & \textbf{Flags} \\
  \midrule
  \endhead
  \bottomrule
  \endfoot
  \code{--enable-openacc} & \code{-acc=gpu -cuda -gpu=ccnative,lineinfo
    -Minfo=accel} \\
  \code{--enable-openacc-debug} & \code{-acc=gpu -cuda
    -gpu=ccnative,autocompare -Minfo=accel}, on the \code{NOOPTFC} flags
    for every file \\
  \code{--enable-openacc-managed} & \code{-acc=gpu -cuda
    -gpu=ccnative,lineinfo,mem:managed -Minfo=accel} \\
  \code{--enable-openacc-stdpar} & \code{-stdpar=gpu -cuda
    -gpu=ccnative,lineinfo,mem:unified -Minfo=accel} \\
\end{longtable}

\section{Substituted variables}

Besides the standard autoconf ones, \code{configure} substitutes:

\begin{longtable}{L{4cm} L{9.5cm}}
  \toprule
  \textbf{Variable} & \textbf{Content} \\
  \midrule
  \endhead
  \bottomrule
  \endfoot
  \code{MPIFC} & The MPI Fortran compiler (\code{ACX\_MPI}), used by most
    compile rules and as \code{FCLD}. \\
  \code{FCFLAGS}, \code{F90FLAGS}, \code{NOOPTFC}, \code{OMPFLAGS} & The
    Fortran flags (Section~\ref{sec:fcflags}); \code{F90FLAGS} is a copy
    of \code{FCFLAGS}. \\
  \code{AM\_CPPFLAGS} & The \code{-D} definitions and the include paths of
    \code{external} and NetCDF. \\
  \code{CPPFLAGS} & \code{AM\_CPPFLAGS} plus \code{-I} of \code{Share}. \\
  \code{AM\_LDFLAGS}, \code{LDFLAGS} & Library paths (\code{external},
    NetCDF) and the GPU link flags. \\
  \code{LIBS}, \code{SYSLIBS} & \code{-lrcmlib} of \code{Share}, then
    NetCDF, OASIS and pthread libraries. \\
  \code{NC\_CC}, \code{NC\_PREFIX}, \code{NF\_PREFIX}, \code{NC\_INCLUDES},
    \code{NC\_LDFLAGS}, \code{NC\_LIBS} & The NetCDF installation. \\
  \code{REGCM\_ROOT} & The source directory, also the default install
    prefix. \\
\end{longtable}

\section{Testing a change}

A change to \code{configure.ac} or \code{acinclude.m4} that should not
change the flags is checked by comparing the variables of the generated
top-level \code{Makefile} before and after, for several option sets:

\begin{lstlisting}
./bootstrap.sh
./configure --enable-clm45 --enable-debug
grep -E '^[A-Za-z0-9_]+ = ' Makefile | sort > after.txt
\end{lstlisting}

The values in \code{config.status} are split over several lines and do not
suit a plain comparison. A compiler that is not installed can be imitated
by a small script that prints its banner for \code{--version} and passes
the other calls to an installed compiler, removing the flags that compiler
does not accept; pass the same script as \code{FC} and \code{MPIFC}. A
change that should alter the flags is finished with a clean full build
(\code{make clean; make}) and the model regression tests.

\fussy
